% Section declaración - DOCENCIA E INVESTIGACIÓN
\par{
	
	Economista con sólida formación metodológica y experiencia docente en \textbf{Metodología de la Investigación} a nivel de pregrado en ciencias económicas y sociales. Especializado en acompañamiento integral de estudiantes en el proceso investigativo: desde el planteamiento del problema hasta la sustentación de tesis, con dominio de diseños cuantitativos, cualitativos y mixtos.
	
	Experto en elaboración y validación de instrumentos de investigación, análisis de datos con software estadístico (R, Stata, SPSS), redacción científica bajo normas APA 7 y Vancouver, y formación en ética investigativa. Apasionado por desarrollar competencias investigativas rigurosas en estudiantes y promover la investigación aplicada que genere impacto social.
	
	Busco oportunidades en docencia universitaria o institutos donde pueda contribuir con mi experiencia en tutoría de tesis, dictado de cursos metodológicos, talleres de IA aplicada a la investigación, y promoción de la cultura investigativa, fortaleciendo así el desarrollo académico y profesional de los estudiantes.

	%%%%%%%%%%%%%%%%%%%%%%%%%%%%%%%%%%%%%%%%
	% VERSIÓN ALTERNATIVA - Más enfocada en investigación científica
	%%%%%%%%%%%%%%%%%%%%%%%%%%%%%%%%%%%%%%%%
	%Docente investigador en ciencias económicas y sociales con experiencia en formación metodológica a nivel de pregrado. Especializado en diseño y ejecución de proyectos de investigación cuantitativos y cualitativos, con énfasis en análisis estadístico y econométrico aplicado.
	
	%Combino formación económica con dominio de herramientas de análisis de datos (R, Python, Stata) para desarrollar investigaciones rigurosas y formar investigadores competentes. Experiencia en publicación científica, revisión por pares, y gestión de proyectos de investigación interdisciplinarios.
	
	%Busco oportunidades en instituciones académicas comprometidas con la investigación aplicada y la formación de investigadores éticos y competentes, donde pueda liderar proyectos, dirigir tesis y contribuir a la generación de conocimiento científico de calidad.

	%%%%%%%%%%%%%%%%%%%%%%%%%%%%%%%%%%%%%%%%
	% VERSIÓN ALTERNATIVA - Énfasis en educación superior
	%%%%%%%%%%%%%%%%%%%%%%%%%%%%%%%%%%%%%%%%
	%Docente universitario con experiencia en cursos de investigación, estadística y economía aplicada. Comprometido con la innovación pedagógica y el uso de herramientas digitales para el aprendizaje significativo. Especializado en metodología de la investigación, análisis cuantitativo y formación de competencias investigativas.
	
	%Experto en diseño curricular, evaluación por competencias y uso de plataformas educativas (Moodle, Google Classroom, Canvas). Capaz de impartir cursos en modalidad presencial, virtual e híbrida, adaptándome a las necesidades del entorno educativo actual.
	
	%Busco contribuir al desarrollo académico de instituciones educativas de nivel superior, aportando mi experiencia en docencia, tutoría y gestión académica para la formación integral de profesionales competentes y éticos.

	%%%%%%%%%%%%%%%%%%%%%%%%%%%%%%%%%%%%%%%%
	% VERSIÓN ALTERNATIVA - Coordinación académica
	%%%%%%%%%%%%%%%%%%%%%%%%%%%%%%%%%%%%%%%%
	%Economista con experiencia en docencia universitaria y coordinación de programas académicos. Especializado en metodología de investigación, análisis de datos y diseño curricular. Capacidad demostrada en gestión de equipos docentes, seguimiento de indicadores académicos y mejora continua de programas educativos.
	
	%Combino experiencia docente con habilidades de gestión académica: planificación de syllabi, coordinación de prácticas preprofesionales, evaluación de calidad educativa y vinculación con sectores productivos. Orientado a resultados y comprometido con la excelencia académica.
	
	%Busco oportunidades en coordinación académica, dirección de programas de pregrado o posgrado, o jefatura de áreas de investigación donde pueda aplicar mi visión integral de la educación superior y mi capacidad de liderazgo pedagógico.
	
}
